
\section*{5. Separation of groups and homogeneous spaces}
\phantomsection\label{sec:sga3-VIB-5}%
\addcontentsline{toc}{section}{5. Separation of groups and homogeneous spaces}

\medskip
\phantomsection\label{prop:sga3-VIB-5.1}%
\noindent\textbf{Proposition 5.1.---}
\emph{For an \(S\)-group \(G\) to be separated, it is necessary
and sufficient that the identity section of \(G\) be a closed
immersion.}

The condition is necessary (EGA~I, 5.4.6); it is sufficient by
virtue of the following cartesian diagram (cf. VI\(_A\), 0.3):
\begin{center}
  \phantomsection\label{diag:sga3-VIB-5.1-cartesian}%
  \begin{tikzcd}[row sep=large,column sep=huge]
    G\times_S G
      \arrow[r,"{\mu\circ(\operatorname{id}_G\times c)}"] &
    G \\
    G
      \arrow[u,"\Delta_{G/S}"]
      \arrow[r,"\pi"] &
    S
      \arrow[u,"\varepsilon"']
  \end{tikzcd}\,.
\end{center}

\medskip
\phantomsection\label{prop:sga3-VIB-5.2}%
\noindent\textbf{Proposition 5.2.---}
\emph{If \(S\) is discrete, every \(S\)-group is separated.}

Indeed,
\[
  S=\coprod_{s\in S}\operatorname{Spec}\mathcal O_{S,s},
\]
and, by EGA~I, 5.5.5, it is enough to show that, for every
\(s\in S\),
\[
  G\times_S\operatorname{Spec}\mathcal O_{S,s}
\]
is separated. This follows from VI\(_A\), \SGARefLink{sga3:VIA:0.3}{0.3}, since
\(\mathcal O_{S,s}\) is a local ring of dimension \(0\).%

\medskip
\phantomsection\label{thm:sga3-VIB-5.3}%
\noindent\textbf{Theorem 5.3.---}
\emph{Let \(S\) be a scheme, let \(G\) be an \(S\)-group scheme
locally of finite presentation over \(S\) and universally open over
\(S\) in a neighborhood of the identity section, and let \(X\) be
an \(S\)-scheme on which \(G\) acts in such a way that the morphism}
\[
  \Phi:G\times_S X\longrightarrow X\times_S X,
  \qquad (g,x)\longmapsto(gx,x),
\]
\emph{is surjective. Suppose moreover that, for every \(s\in S\):}
\begin{enumerate}
\item[\textup{(i)}] \emph{there exists an open subscheme \(U\) of
  \(X\), separated over \(S\), such that \(U_s\) is dense in
  \(X_s\);}
\item[\textup{\SGAInlineRefTarget{thm:sga3-VIB-5.3:item:ii}(ii)}] \emph{the fiber \(X_s\) is locally of finite
  type over \(\kappa(s)\).}%
\footnote{\label{note:sga3-VIB-34}Editor's note: we have corrected
the original by assuming that \(G\) is universally open over \(S\)
in a neighborhood of the identity section and by adding hypothesis
\textup{\SGARefLink{thm:sga3-VIB-5.3:item:ii}{(ii)}}; see below for examples, due to O.\ Gabber, showing
that the statements \hyperref[thm:sga3-VIB-5.3]{5.3} and \hyperref[cor:sga3-VIB-5.4]{5.4} of the original are false without
additional hypotheses. We also point out that \hyperref[thm:sga3-VIB-5.3]{Theorem 5.3} is a
reworked version of \hyperref[thm:sga3-VIB-5.3A]{Theorem 5.3A} below, which appears in the 1965
edition of SGAD and is due to M.\ Raynaud; compare the notes
\((*)\) in \SGARefLink{sga3:x:lemma:8-5}{Expos\'e~X, 8.5} and 8.8.

\smallskip
\phantomsection\label{thm:sga3-VIB-5.3A}%
\noindent\textbf{Theorem 5.3A (Raynaud).---}
\emph{Let \(G\) be an \(S\)-group locally of finite presentation,
universally open over \(S\), and with connected fibers. Then \(G\)
is separated over \(S\). More generally, every \(S\)-scheme \(X\)
locally of finite presentation over \(S\), equipped with an action
of \(G\) such that the morphism
\(\Phi:G\times_SX\to X\times_SX\),
\((g,x)\mapsto(gx,x)\), is surjective, is separated over \(S\).}}
\end{enumerate}
\emph{Then \(X\) is separated over \(S\).}

\medskip
\phantomsection\label{cor:sga3-VIB-5.4}%
\noindent\textbf{Corollary 5.4.---}
\emph{Let \(S,G,X\) be as in the preliminary hypotheses of \SGARefLink{thm:sga3-VIB-5.3}{5.3}.
Suppose moreover that \(X\) has fibers locally of finite type and
connected. Then:}
\begin{enumerate}
\item[\textup{\SGAInlineRefTarget{cor:sga3-VIB-5.4:item:i}(i)}] \emph{\(X\) is separated over \(S\).}
\item[\textup{(ii)}] \emph{If there exists an open subset \(V\) of
  \(X\), quasi-compact over \(S\), which meets every nonempty fiber
  \(X_s\),}%
\emph{then \(X\) is quasi-compact over \(S\).}
\end{enumerate}

\emph{Proof.}
\textup{\SGARefLink{cor:sga3-VIB-5.4:item:i}{(i)}} Let \(s\in S\) be such that \(X_s\ne\varnothing\).
The morphism
\[
  \Phi_s:
  G_s\times_{\kappa(s)}X_s
  \longrightarrow
  X_s\times_{\kappa(s)}X_s
\]
obtained from \(\Phi\) by base change is surjective. Since \(X_s\)
is connected, it is irreducible by VI\(_A\), \SGARefLink{prop:sga3-VIA-2.5.4}{2.5.4}. Thus, if \(U\)
is an affine open subset of \(X\) such that \(U_s\) is nonempty,
then \(U_s\) is dense in \(X_s\), and the theorem applies.

To prove \textup{\SGARefLink{concl:sga3-VIA-2.5.4-ii}{(ii)}}, we may suppose that \(S\) is affine. Then
\(V\) is quasi-compact and, by \hyperref[prop:sga3-VIB-3.5]{3.5}, there exists a quasi-compact
open subset \(U\) of \(G\) containing \(\underline{G^0}\). Let
\(s\in S\) be such that \(X_s\ne\varnothing\). Then \(X_s\) is
irreducible (VI\(_A\), \SGARefLink{prop:sga3-VIA-2.6.6}{2.6.6}), and, since \(U_s\) contains
\(\underline{G_s^0}\), the morphism
\[
  U_s\times_{\kappa(s)}V_s\longrightarrow X_s,
  \qquad (g,x)\longmapsto gx,
\]
is surjective (VI\(_A\), \SGARefLink{prop:sga3-VIA-2.6.4}{2.6.4}). Consequently the morphism
\[
  U\times_SV\longrightarrow X
\]
is surjective, and \(X\) is therefore quasi-compact, since \(U\)
and \(V\) are. As \(S\) was assumed affine, hence separated, it
follows that \(X\) is quasi-compact over \(S\)
(cf. EGA~I, 6.6.4\textup{(v)}).

\medskip
\phantomsection\label{rem:sga3-VIB-5.4.1}%
\noindent\textbf{Remark 5.4.1.---}%
\footnote{\label{note:sga3-VIB-36}Editor's note: we have added this
remark.}
We shall see in the course of the proof that the conclusion of
\hyperref[thm:sga3-VIB-5.3]{Theorem 5.3} remains true if one assumes only:

\begin{quote}
\textup{(i')} there exists an open subscheme \(U\) of \(X\),
separated over \(S\), such that \(U_s\) is dense in every
irreducible component of \(X_s\).
\end{quote}

This follows from \textup{\SGARefLink{cor:sga3-VIB-5.4:item:i}{(i)}} if \(X_s\) locally has a finite
number of irreducible components, as is the case under hypothesis
\textup{(ii)}. We shall also see later that \SGARefLink{cor:sga3-VIB-5.4}{5.4} remains true under
the hypothesis that every fiber \(X_s\) is quasi-separated and
connected.
